\chapter{Electric Power and Energy}\label{ch:g9:electric-power-energy}

The electricity bill does not charge for \omterm{def:g8:voltage-voltmeter:voltage}{volts}, nor for
\omterm{def:g8:intensity-ammeter:intensity}{amperes} --- it charges for something your family's appliances
have been quietly multiplying together all month. This chapter
weds the two electrical quantities into a third, gives fast
and slow spending their proper names, and ends at the meter by
the front door: physics you can audit against real money.

\section{Power: the rate of delivery}

\begin{definition}[Power and the watt]\label{def:g9:electric-power-energy:power}
The \emph{power}\index{power} $P$ of a device is the rate at
which it converts \omterm{def:g5:everyday-energy:energy}{energy} --- \omterm{def:g9:kinetic-energy-safety:joule}{joules} per second. Its \omterm{def:g6:measurement-in-science:measuring}{unit} is
the \emph{watt}\index{watt} (\unit{W}): one \omterm{def:g9:kinetic-energy-safety:joule}{joule} each
second, honoring the engineer of the steam age; the
\emph{kilowatt} (\qty{1}{kW} $=$ \qty{1000}{W}) serves the
kitchen's heavyweights. Power is not an amount of \omterm{def:g5:everyday-energy:energy}{energy} but
a \emph{pace} of spending --- a tap's flow, not a bucket's
content.
\end{definition}

\begin{proposition}[Electric power]\label{prop:g9:electric-power-energy:UI}
A device fed \omterm{def:g8:voltage-voltmeter:voltage}{voltage} $U$ and drawing current $I$ receives
electric \omterm{def:g9:electric-power-energy:power}{power}
\[
P = U \times I ,
\]
\omterm{def:g9:electric-power-energy:power}{watts} from \omterm{def:g8:voltage-voltmeter:voltage}{volts} times \omterm{def:g8:intensity-ammeter:intensity}{amperes}. The formula computes in all
three directions, and its second recipe runs the \omterm{def:g7:short-circuits-safety:fuse}{fuse} box:
$I = P/U$ tells what current an appliance's rated \omterm{def:g9:electric-power-energy:power}{power} will
pull from the mains.
\end{proposition}

\begin{example}[The household ladder]\label{ex:g9:electric-power-energy:ladder}
Read the plates on the family's machines: an indicator \omterm{def:g1:light-and-shadows:source}{light},
\qty{1}{W}; a modern lamp, \qty{5}{W} (its glowing ancestor
spent sixty for the same \omterm{def:g1:light-and-shadows:source}{light} --- the difference was \omterm{def:g5:heat-and-insulation:heat}{heat});
a laptop, \qty{50}{W}; a television, \qty{100}{W}; a
washing machine heating its water, \qty{2000}{W}; the
kettle, \qty{2300}{W}; an instant water heater, up to
\qty{9000}{W}. The pattern of the ladder: anything whose job
is \emph{heating} climbs to the kilowatts --- warmth is the
costliest form of \omterm{def:g5:everyday-energy:energy}{energy} to supply.
\end{example}

\begin{example}[The fuse box arithmetic]\label{ex:g9:electric-power-energy:fuses}
The kettle on the \qty{230}{V} mains: $I = 2300 \div 230 =
\qty{10}{A}$ --- the \omterm{def:g8:intensity-ammeter:intensity}{ten-ampere} appliance of the old
chapters, now derived. The \qty{9000}{W} water heater: nearly
\qty{40}{A}, demanding its own stout \omterm{def:g5:series-parallel-circuits:parallel}{branch} and breaker. And
the overloaded multi-socket of the fire investigation:
kettle, heater and fryer summing their $P/U$ currents past
the \omterm{def:g5:series-parallel-circuits:parallel}{branch}'s rating --- three innocents, one arithmetic
crime, exactly as charged.
\end{example}

\section{Energy: the amount delivered}

\begin{proposition}[Energy from power and time]\label{prop:g9:electric-power-energy:Pt}
A device of \omterm{def:g9:electric-power-energy:power}{power} $P$ running for a time $t$ converts the
\omterm{def:g5:everyday-energy:energy}{energy}
\[
E = P \times t ,
\]
\omterm{def:g9:kinetic-energy-safety:joule}{joules} from \omterm{def:g9:electric-power-energy:power}{watts} times seconds. Fast tap or long trickle,
the bucket fills by the product: a \qty{2300}{W} kettle's
three minutes ($2300 \times 180 \approx \qty{4.1e5}{J}$)
spends more than a \qty{5}{W} lamp's whole day ($5 \times
86400 = \qty{4.3e5}{J}$) --- almost exactly a tie, in fact:
one boiling matches one day of good \omterm{def:g1:light-and-shadows:source}{light}.
\end{proposition}

\begin{definition}[The kilowatt-hour]\label{def:g9:electric-power-energy:kwh}
Household energies make ungainly joule-counts, so the meter
speaks \emph{kilowatt-hours}\index{kilowatt-hour}
(\unit{kWh}): the \omterm{def:g5:everyday-energy:energy}{energy} of one kilowatt sustained for one
hour,
\[
\qty{1}{kWh} = 1000 \times 3600 = \qty{3.6e6}{J}
\]
--- three point six million \omterm{def:g9:kinetic-energy-safety:joule}{joules} per \omterm{def:g6:measurement-in-science:measuring}{unit} on the bill.
With $P$ in kilowatts and $t$ in \omterm{ex:g2:measuring-time:clock}{hours}, $E = P \times t$
delivers \omterm{ex:g2:measuring-time:clock}{kilowatt-hours} directly: the formula the meter
lives by.
\end{definition}

\begin{method}[Auditing an appliance]\label{met:g9:electric-power-energy:audit}
For any machine and month:
\begin{enumerate}
  \item read its \omterm{def:g9:electric-power-energy:power}{power} $P$ on the plate (or compute
        $U \times I$);
  \item honestly estimate its running time $t$ --- per day,
        then per month;
  \item $E = P \times t$ in \omterm{def:g9:electric-power-energy:kwh}{kilowatt-hours} ($P$ in
        \unit{kW}, $t$ in \omterm{ex:g2:measuring-time:clock}{hours});
  \item multiply by the tariff (take about $0.25$ currency
        \omterm{def:g6:measurement-in-science:measuring}{units} per \unit{kWh}) for the month's cost.
\end{enumerate}
The kettle, six minutes daily: $2.3 \times 0.1 \times 30 =
\qty{6.9}{kWh}$ a month. The television, three \omterm{ex:g2:measuring-time:clock}{hours} daily:
$0.1 \times 3 \times 30 = \qty{9}{kWh}$. Small \omterm{def:g9:electric-power-energy:power}{powers} with
long \omterm{ex:g2:measuring-time:clock}{hours} rival big \omterm{def:g9:electric-power-energy:power}{powers} with short ones --- the audit's
recurring lesson.
\end{method}

\begin{omfigure}
\begin{tikzpicture}
\begin{axis}[omaxis, width=10.6cm, height=5.4cm,
    xlabel={time (\unit{h})}, ylabel={power (\unit{kW})},
    xtick={0,1,2,3}, ytick={0,1,2,2.3},
    ymin=0, ymax=2.9, xmin=0, xmax=3.4]
  \addplot[very thick, omThm] coordinates
    {(0,2.3) (0.1,2.3)};
  \addplot[omThm, fill=omThm, fill opacity=0.25, draw=none]
    coordinates {(0,0) (0,2.3) (0.1,2.3) (0.1,0)} \closedcycle;
  \addplot[very thick, omDef] coordinates {(0,0.1) (3,0.1)};
  \addplot[omDef, fill=omDef, fill opacity=0.25, draw=none]
    coordinates {(0,0) (0,0.1) (3,0.1) (3,0)} \closedcycle;
  \node[right, font=\footnotesize, omThm] at (axis cs:0.15,2.0)
    {kettle: tall and brief};
  \node[above, font=\footnotesize, omDef] at (axis cs:1.9,0.14)
    {television: low and long};
\end{axis}
\end{tikzpicture}

{\small \omterm{def:g5:everyday-energy:energy}{Energy} as area: \omterm{def:g9:electric-power-energy:power}{power} times time. The kettle's tall
sliver and the television evening's long strip can enclose
comparable areas --- comparable \omterm{def:g9:electric-power-energy:kwh}{kilowatt-hours}.}
\end{omfigure}

\section{The meter and the bill}

\begin{example}[Reading the bill]\label{ex:g9:electric-power-energy:bill}
The meter by the front door counts every \omterm{def:g9:electric-power-energy:kwh}{kilowatt-hour} that
enters the house; the bill is its month's difference times
the tariff. A typical family's $250$ monthly \omterm{def:g6:measurement-in-science:measuring}{units} decompose
by audit: water heating and radiators first (the kilowatt
club), then the cold appliances --- modest \omterm{def:g9:electric-power-energy:power}{watts}, but
running always --- then cooking, washing, \omterm{def:g1:light-and-shadows:source}{light} and
electronics. The audit's \omterm{def:g9:electric-power-energy:power}{power}: it finds the real \omterm{def:g3:levers-and-scales:lever}{levers}.
Replacing ten old \qty{60}{W} \omterm{def:g3:first-electric-circuit:bulb}{bulbs} with \qty{5}{W} ones
saves more than unplugging every charger in the house a
thousand times over.
\end{example}

\begin{example}[The standby vampires]\label{ex:g9:electric-power-energy:standby}
Small print, long \omterm{ex:g2:measuring-time:clock}{hours}: a television's standby \omterm{def:g1:light-and-shadows:source}{light} sips
\qty{0.5}{W} --- but for all $720$ \omterm{ex:g2:measuring-time:clock}{hours} of the month:
$0.0005 \times 720 = \qty{0.36}{kWh}$. One vampire is
harmless; a house with twenty sippers at a \omterm{def:g9:electric-power-energy:power}{watt} each pays
$0.02 \times 720 \approx \qty{14}{kWh}$ monthly for
nothing. \omterm{def:g9:electric-power-energy:power}{Power} times time forgives no term: the second
factor bites when the first one hides.
\end{example}

\begin{remark}[Power in the wider world]\label{rem:g9:electric-power-energy:wider}
The \omterm{def:g9:electric-power-energy:power}{watt}, born electric here, \omterm{def:g6:measurement-in-science:measuring}{measures} every \omterm{def:g5:everyday-energy:energy}{energy} pace.
Your resting body idles near \qty{100}{W} --- a glowing
\omterm{def:g3:first-electric-circuit:bulb}{ancestor-bulb} of warmth (recall the crowded classroom's
stuffiness: thirty pupils, three kilowatts of biology). A
cyclist sustains \qty{200}{W}; a car's engine unleashes
tens of kilowatts; the town's power station, hundreds of
millions of \omterm{def:g9:electric-power-energy:power}{watts}; the sunlight falling on the town,
comfortably more --- the ceiling under which all the other
paces work. Next chapter follows the grid's \omterm{def:g9:electric-power-energy:power}{watts} to their
spinning source.
\end{remark}

\section{Exercises}

\begin{exercise}[$\star$]\label{exo:g9:electric-power-energy:1}
Define \omterm{def:g9:electric-power-energy:power}{power} and its \omterm{def:g6:measurement-in-science:measuring}{unit} --- tap or bucket? Write both
formulas of the chapter.
\end{exercise}

\begin{exercise}[$\star$]\label{exo:g9:electric-power-energy:2}
Compute \omterm{def:g9:electric-power-energy:power}{powers}: a lamp at \qty{230}{V} drawing
\qty{0.022}{A}; a starter \omterm{ex:g6:circuits-and-safety:motor}{motor} at \qty{12}{V} drawing
\qty{100}{A}.
\end{exercise}

\begin{exercise}[$\star$]\label{exo:g9:electric-power-energy:3}
The fuse-box recipe: what currents do a \qty{2300}{W}
kettle and a \qty{700}{W} microwave pull at \qty{230}{V}?
\end{exercise}

\begin{exercise}[$\star$]\label{exo:g9:electric-power-energy:4}
Why do heating appliances crowd the top of the household
\omterm{def:g9:electric-power-energy:power}{power} ladder?
\end{exercise}

\begin{exercise}[$\star$]\label{exo:g9:electric-power-energy:5}
Convert: \qty{1}{kWh} to \omterm{def:g9:kinetic-energy-safety:joule}{joules} (show the two factors).
Why did the meter's makers abandon the \omterm{def:g9:kinetic-energy-safety:joule}{joule}?
\end{exercise}

\begin{exercise}[$\star$]\label{exo:g9:electric-power-energy:6}
\omterm{def:g5:everyday-energy:energy}{Energy} check: which spends more, a \qty{2}{kW} heater for
half an hour or a \qty{100}{W} television for a whole
evening of five \omterm{ex:g2:measuring-time:clock}{hours}?
\end{exercise}

\begin{exercise}[$\star$]\label{exo:g9:electric-power-energy:7}
Audit one machine: a \qty{50}{W} laptop used four \omterm{ex:g2:measuring-time:clock}{hours}
daily, for a thirty-day month, at $0.25$ per \unit{kWh}.
\omterm{def:g5:everyday-energy:energy}{Energy} and cost?
\end{exercise}

\begin{exercise}[$\star$]\label{exo:g9:electric-power-energy:8}
What does the front-door meter count, and how does the
bill turn its count into money?
\end{exercise}

\begin{exercise}[$\star\star$]\label{exo:g9:electric-power-energy:9}
The \omterm{def:g3:first-electric-circuit:bulb}{bulb} swap: ten \qty{60}{W} ancestors versus ten
\qty{5}{W} moderns, four \omterm{ex:g2:measuring-time:clock}{hours} nightly, thirty nights.
Compute both monthly energies and the saving --- in
\omterm{def:g9:electric-power-energy:kwh}{kilowatt-hours} and at $0.25$ each.
\end{exercise}

\begin{exercise}[$\star\star$]\label{exo:g9:electric-power-energy:10}
A hair dryer's plate reads ``230 V, 2000 W''. Its current?
May it share a \qty{16}{A} \omterm{def:g5:series-parallel-circuits:parallel}{branch} with the \qty{2300}{W}
kettle, both running? Show the sum.
\end{exercise}

\begin{exercise}[$\star\star$]\label{exo:g9:electric-power-energy:11}
An electric car's \omterm{def:g3:first-electric-circuit:battery}{battery} stores \qty{60}{kWh}. How long
could that run the \qty{2.3}{kW} kettle? The \qty{5}{W}
lamp? And what does the pair of answers teach about the
size of transport's \omterm{def:g5:everyday-energy:energy}{energy} appetite?
\end{exercise}

\begin{exercise}[$\star\star\star$]\label{exo:g9:electric-power-energy:12}
The kettle test, end to end: \qty{1.5}{L} of water from
\qty{20}{\celsius} to the boil needs about \qty{5.0e5}{J}
of \omterm{def:g5:heat-and-insulation:heat}{heat}. Time the family kettle (\qty{2300}{W}): predict
its boiling time from $E = P t$, then explain why the real
kettle takes a little longer --- where do the missing
\omterm{def:g9:kinetic-energy-safety:joule}{joules} go, by the oldest law of the \omterm{def:g5:everyday-energy:energy}{energy} chapters?
\end{exercise}

\section{Problem: The Home Energy Audit}

\begin{problem}[{Weekend problem --- the family energy
audit; plates, meters and the month's bill; finding the
real levers}]\label{pb:g9:electric-power-energy:1}
Armed with the two formulas and the tariff ($0.25$ per
\unit{kWh}), the family audits a month (30 days). The
inventory: kettle \qty{2300}{W}, $6$ minutes daily;
television \qty{100}{W}, $4$ \omterm{ex:g2:measuring-time:clock}{hours} daily; refrigerator ---
special case below; ten lamps \qty{5}{W} each, $5$ \omterm{ex:g2:measuring-time:clock}{hours}
daily; washing machine \qty{2}{kW}, one $1$-hour hot cycle
every other day; water heater \qty{2.5}{kW}, $2$ \omterm{ex:g2:measuring-time:clock}{hours}
daily; standby sippers totaling \qty{10}{W}, always on.

\medskip\noindent\textbf{Part I --- Machine by machine.}
\begin{enumerate}
  \item The kettle's month, in \unit{kWh}.
  \item The television's and the ten lamps' months.
  \item The washing machine's ($15$ cycles) and the water
        heater's months.
  \item The standby sippers' month --- the second factor at
        work.
\end{enumerate}

\medskip\noindent\textbf{Part II --- The refrigerator and
the sum.}
The refrigerator's compressor runs at \qty{120}{W}, but
only a third of the time, cycling day and night.
\begin{enumerate}[resume]
  \item Its effective average \omterm{def:g9:electric-power-energy:power}{power}, and its month in
        \unit{kWh}.
  \item Total the audit's seven lines. Which three lines
        dominate?
  \item The month's bill at the tariff --- and the water
        heater's share of it in currency.
  \item The meter read $2150$ \omterm{def:g6:measurement-in-science:measuring}{units} at the month's start.
        Predict its end-of-month reading.
\end{enumerate}

\medskip\noindent\textbf{Part III --- The \omterm{def:g3:levers-and-scales:lever}{levers}.}
\begin{enumerate}[resume]
  \item Proposal one: shorter showers --- the heater drops
        to \qty{1.5}{h} daily. Monthly saving, \omterm{def:g5:everyday-energy:energy}{energy} and
        money?
  \item Proposal two: a public campaign urges unplugging
        phone chargers (\qty{0.2}{W} each when idle).
        Audit one charger's idle month and compare with
        proposal one; what does the contrast teach about
        audits versus slogans?
  \item Proposal three: grandmother suggests the old habit
        of filling the kettle only as needed --- half the
        water, half the \omterm{def:g5:everyday-energy:energy}{energy}. New kettle line, and the
        month's saving?
  \item Write the audit's conclusions: three sentences ---
        where the \omterm{def:g9:electric-power-energy:kwh}{kilowatt-hours} truly live, which formula
        factor the standby case teaches, and which single
        household change pays best.
\end{enumerate}
\end{problem}
